\documentclass[10pt]{article}
\usepackage[margin=0.76in]{geometry}
\usepackage{amsmath,amssymb,amsthm,mathtools}
\usepackage{graphicx}
\usepackage{microtype}
\usepackage[hidelinks]{hyperref}

\newtheorem{theorem}{Theorem}
\newtheorem{corollary}[theorem]{Corollary}
\newtheorem{lemma}[theorem]{Lemma}
\newcommand{\Tr}{\operatorname{Tr}}
\newcommand{\Fh}{\widehat F}
\newcommand{\Dh}{\widehat D}
\newcommand{\eps}{\varepsilon}

\title{Optimal precision dependence for geometric Rényi quasi-entropy\\
in the previously open range $1<\alpha\leq2$}
\author{Candidate full solution for \texttt{fa\_banach\_001}}
\date{17 August 2026}

\begin{document}
\maketitle

\paragraph{Status.}
Candidate full solution, likely valid; human review recommended.  For every
fixed $\alpha\in(1,2]$, the source algorithm's
$\widetilde O_\alpha(1/\eps)$ precision dependence is optimal, even on
commuting full-rank qubit states with constant condition bounds.  The same
example also gives matching precision lower bounds for the logarithmic
geometric Rényi entropy and a quadratic copy lower bound.

\section{Source problem}

Liu, Wang, Wilde, and Zhang define the geometric Rényi relative quasi-entropy
of positive definite density matrices by
\begin{equation}\label{eq:def}
 \Fh_\alpha(\rho,\sigma)
 =\Tr(\sigma\#_\alpha\rho)
 =\Tr\!\left[\sigma
   (\sigma^{-1/2}\rho\sigma^{-1/2})^\alpha\right],
 \qquad \alpha\in(0,1)\cup(1,2].
\end{equation}
They construct an additive-$\eps$ estimator using
$\widetilde O_\alpha(1/\eps)$ purification-oracle queries when the lower
spectral bounds are fixed.  Their Lemma~20 proves a matching
$\Omega_\alpha(1/\eps)$ lower bound for $0<\alpha<1$.  Immediately afterward
they ask whether optimality still holds for $1<\alpha\leq2$
\cite[PDF p.~24, after Lemma~20]{source}.

\begin{figure}[ht]
 \centering
 \includegraphics[width=\textwidth]{figures/open_problem_crop.png}
 \caption{The source's lower bound for $0<\alpha<1$ and the explicit open
 problem for $1<\alpha\leq2$ (arXiv:2405.00673, PDF p.~24).}
\end{figure}

\paragraph{Transcription.}
``It remains an open problem [whether] optimality still holds if
$\alpha\in(1,2]$.''  The source notes that $\Fh_\alpha\geq1$ in this range
and anticipates that different techniques may be required.

\section{Idea of the proof}

No new noncommutative construction is needed.  Use the source's own commuting
Bernoulli hard pair and inspect the scalar target as the Bernoulli bias varies.
The first derivative is nonzero on both sides of $\alpha=1$.  It changes sign
when one keeps the source's ordering of the two hard instances, but query
discrimination only uses the absolute output gap.  Reversing the decision
threshold preserves a linear gap while the purification oracles remain only
$O(t)$ apart in Hellinger distance.  Belovs' oracle-discrimination theorem then
forces $\Omega(1/t)$ queries.  Taking $t\asymp_\alpha\eps$ proves the result.

\section{Query lower bound}

For $t\in[0,1/4]$, put
\begin{equation}\label{eq:states}
 \rho_t=\begin{pmatrix}(1+t)/2&0\\0&(1-t)/2\end{pmatrix},
 \qquad
 \eta=\begin{pmatrix}1/4&0\\0&3/4\end{pmatrix}.
\end{equation}
Both $\rho_t$ and $\eta$ dominate $I/4$.  Since they commute,
\begin{align}
 g_\alpha(t)
 &:=\Fh_\alpha(\rho_t,\eta) \label{eq:g}\\
 &=\left(\frac14\right)^{1-\alpha}
       \left(\frac{1+t}{2}\right)^\alpha
   +\left(\frac34\right)^{1-\alpha}
       \left(\frac{1-t}{2}\right)^\alpha .\notag
\end{align}

\begin{lemma}[Linear separation]\label{lem:gap}
For every fixed $\alpha>1$, there are constants $c_\alpha,t_\alpha>0$
such that
\[
 g_\alpha(2t)-g_\alpha(t)\geq c_\alpha t
 \qquad(0<t\leq t_\alpha).
\]
The same assertion holds with $g_\alpha$ replaced by
$r_\alpha=(\alpha-1)^{-1}\log g_\alpha$.
\end{lemma}

\begin{proof}
Differentiation at zero gives
\begin{equation}\label{eq:derivative}
 g_\alpha'(0)
 =\alpha2^{-\alpha}\left[
       \left(\frac14\right)^{1-\alpha}
      -\left(\frac34\right)^{1-\alpha}\right]>0,
\end{equation}
because $1-\alpha<0$.  By continuity, $g_\alpha'(u)$ is bounded below by
$g_\alpha'(0)/2$ on a sufficiently short interval $[0,2t_\alpha]$.
Integrating over $[t,2t]$ proves the first assertion.  Moreover,
\[
 r_\alpha'(0)=\frac{g_\alpha'(0)}
 { (\alpha-1)g_\alpha(0)}>0,
\]
so the identical continuity argument proves the second.
\end{proof}

We use the following query lower bound in precisely the oracle model used by
the source.  If $p,q$ are classical distributions and unitary oracles prepare
purifications with amplitudes $\sqrt{p_j}$ and $\sqrt{q_j}$, respectively,
then distinguishing the two oracles with constant success probability takes
$\Omega(1/H(p,q))$ queries, where
\[
 H(p,q)^2=\frac12\sum_j(\sqrt{p_j}-\sqrt{q_j})^2.
\]
This is Belovs' distribution-discrimination theorem
\cite[Theorem~4]{belovs}, also quoted as the lower-bound input in the source.

\begin{theorem}[Resolution of the open problem]\label{thm:main}
Fix $\alpha\in(1,2]$.  Any quantum algorithm which, given purification
oracles for $\rho$ and $\sigma$ satisfying $\rho,\sigma\geq I/4$, estimates
$\Fh_\alpha(\rho,\sigma)$ to additive error $\eps$ with success probability
at least $2/3$ requires
\[
 \Omega_\alpha(1/\eps)
\]
queries as $\eps\downarrow0$.
\end{theorem}

\begin{proof}
Let $p_t=((1+t)/2,(1-t)/2)$, the eigenvalue distribution of $\rho_t$.
For $0<t\leq1/4$, every coordinate of both $p_t$ and $p_{2t}$ is at least
$1/4$.  Hence
\begin{align*}
 H(p_t,p_{2t})^2
 &=\frac12\sum_{j=0}^1
   \frac{(p_t(j)-p_{2t}(j))^2}
        {(\sqrt{p_t(j)}+\sqrt{p_{2t}(j)})^2}
 \leq \frac{t^2}{4},
\end{align*}
and in particular $H(p_t,p_{2t})\leq t$.

Let an estimator as in the theorem have additive error $\eps$.  By
Lemma~\ref{lem:gap}, for all sufficiently small $\eps$ we may choose
$t=C_\alpha\eps$ so that $t\leq\min(t_\alpha,1/4)$ and
\[
 g_\alpha(2t)-g_\alpha(t)\geq4\eps.
\]
Give the estimator either a purification oracle for $\rho_t$ or one for
$\rho_{2t}$ as its first input, and give it the fixed purification oracle for
$\eta$ as its second input.  Comparing the output with the midpoint of
$g_\alpha(t)$ and $g_\alpha(2t)$ distinguishes the two cases whenever the
estimation succeeds.  Belovs' theorem therefore forces
\[
 \Omega(1/H(p_t,p_{2t}))=\Omega(1/t)=\Omega_\alpha(1/\eps)
\]
queries to the variable oracle, and hence at least that many total queries.
All input states dominate $I/4$, as required.
\end{proof}

Combining Theorem~\ref{thm:main} with the source upper bound proves that, for
every fixed $\alpha\in(1,2]$, the precision dependence is
$\widetilde\Theta_\alpha(1/\eps)$.  The endpoint $\alpha=2$ is included.
At the excluded point $\alpha=1$, $\Fh_1(\rho,\sigma)=\Tr\rho=1$, so the
change at $\alpha=1$ is genuine rather than a defect of the construction.

\section{Two strengthenings}

The geometric Rényi entropy is
\[
 \Dh_\alpha(\rho\Vert\sigma)
   =\frac{1}{\alpha-1}\log\Fh_\alpha(\rho,\sigma).
\]
The second part of Lemma~\ref{lem:gap}, followed verbatim by the proof of
Theorem~\ref{thm:main}, gives the following.

\begin{corollary}[Entropy precision]
For each fixed $\alpha\in(1,2]$, additive-$\eps$ estimation of
$\Dh_\alpha(\rho\Vert\sigma)$ from purification oracles requires
$\Omega_\alpha(1/\eps)$ queries even for commuting qubit states dominating
$I/4$.  Thus the source's $\widetilde O_\alpha(1/\eps)$ entropy algorithm also
has optimal precision dependence.
\end{corollary}

\begin{corollary}[Copy lower bounds]
For either $\Fh_\alpha$ or $\Dh_\alpha$, estimation to additive error $\eps$
from copies of the input states requires $\Omega_\alpha(1/\eps^2)$ copies,
under the same fixed-$\alpha$ and constant-conditioning assumptions.
\end{corollary}

\begin{proof}
For product distributions,
\[
 H(p_t^{\otimes m},p_{2t}^{\otimes m})^2
 =1-(1-H(p_t,p_{2t})^2)^m
 \leq mH(p_t,p_{2t})^2\leq mt^2/4.
\]
Total variation is at most $\sqrt2$ times Hellinger distance.  Therefore a
constant-bias distinction requires $m=\Omega(1/t^2)$.  The density matrices
commute, so even an arbitrary joint quantum measurement on their tensor
powers has exactly the corresponding classical trace-distance limitation.
Taking $t\asymp_\alpha\eps$ as above proves the claim; copies of the fixed
state $\eta$ convey no information about the case.
\end{proof}

\section{Verification, novelty, and limitations}

The derivative, sign, entropy derivative, finite gaps, and Hellinger estimate
were checked numerically for $\alpha=1.01,1.1,1.25,1.5,1.75,2$ and four scales
$10^{-2}$ through $10^{-5}$ by \texttt{code/check\_hard\_instance.py}; all
checks passed.  This finite computation is only a sanity check, not part of
the proof.

The novelty search covered the run registry and all cheap solution/attempt
indexes; exact searches for the source's open sentence; searches combining
``geometric Rényi relative quasi-entropy'', $\alpha>1$, quantum query
complexity, and precision lower bounds; the current arXiv record; and the 2025
npj Quantum Information publication.  Both current source versions still
state the problem as open, and no later answer was found through 17 August
2026.  The result is elementary and very close to the source proof, so expert
review should focus on the purification-oracle convention and on confirming
that the sign reversal was not handled under different notation elsewhere.

The theorem concerns precision dependence for fixed $\alpha$ and constant
spectral lower bounds.  It does not optimize dependence on $\alpha$, dimension,
rank, or condition numbers, and it does not close the gap between the
$\Omega(1/\eps^2)$ copy lower bound and the source's
$\widetilde O(1/\eps^3)$ copy upper bound.

\begin{thebibliography}{9}
\bibitem{source}
N.~Liu, Q.~Wang, M.~M.~Wilde, and Z.~Zhang,
\emph{Quantum algorithms for matrix geometric means}, arXiv:2405.00673;
npj Quantum Information \textbf{11} (2025), Article 101.

\bibitem{belovs}
A.~Belovs,
\emph{Quantum algorithms for classical probability distributions},
Proceedings of the 27th Annual European Symposium on Algorithms (ESA 2019),
LIPIcs 144, Article 16, 2019.
\end{thebibliography}

\end{document}

